
\chapter{Learning as Function Approximation}

\section{Learning Problems and Data}

Artificial intelligence, in its modern form, is largely concerned with the problem of \emph{learning from data}. At a high level, we are given observations of the world and asked to construct a system that can make predictions, decisions, or generate new data.

We formalize this as follows. Let
\[
\mathcal{X} \subseteq \mathbb{R}^d
\]
denote the \emph{input space} and
\[
\mathcal{Y}
\]
the \emph{output space}. A dataset is a collection of samples
\[
\mathcal{D} = \{(x_i, y_i)\}_{i=1}^n, \quad x_i \in \mathcal{X}, \; y_i \in \mathcal{Y}.
\]

The goal of learning is to construct a function
\[
f : \mathcal{X} \to \mathcal{Y}
\]
that captures the relationship between inputs and outputs.

\medskip

\noindent
\textbf{Examples.}
\begin{itemize}
    \item In \emph{regression}, $\mathcal{Y} \subseteq \mathbb{R}$.
    \item In \emph{classification}, $\mathcal{Y} = \{1, \dots, K\}$.
    \item In \emph{generative modeling}, $\mathcal{Y} = \mathcal{X}$, and the goal is to model the data distribution.
\end{itemize}

\medskip

While this formulation appears simple, it hides the central challenge of machine learning: we only observe finitely many samples, yet we seek a function that performs well on unseen data.

\section{Functions, Models, and Representations}

A learning algorithm does not search over all possible functions. Instead, it selects a function from a \emph{model class} (or \emph{hypothesis space})
\[
\mathcal{H} = \{ f_\theta : \theta \in \Theta \}.
\]

Here:
\begin{itemize}
    \item $\theta$ denotes parameters,
    \item $\Theta$ is the parameter space,
    \item $f_\theta$ is a function determined by $\theta$.
\end{itemize}

\medskip

\noindent
\textbf{Example (Linear model).}
\[
f_\theta(x) = \langle w, x \rangle + b, \quad \theta = (w, b).
\]

\medskip

\noindent
\textbf{Example (Neural network).}
\[
f_\theta(x) = W_L \sigma\big( W_{L-1} \sigma( \cdots \sigma(W_1 x) ) \big),
\]
where $\sigma$ is a nonlinear activation function.

\medskip

The choice of $\mathcal{H}$ encodes assumptions about the structure of the problem. This is known as \emph{inductive bias}, a concept we will revisit in later chapters.

\medskip

\noindent
\textbf{Representation matters.}  
Two models may represent the same class of functions but behave very differently in practice due to:
\begin{itemize}
    \item optimization properties,
    \item numerical stability,
    \item generalization behavior.
\end{itemize}

Thus, a model is not just a function class—it is also a \emph{computational and statistical object}.

\section{Approximation as the Core Paradigm}

At its heart, machine learning is an approximation problem.

\medskip

\noindent
We assume that there exists an unknown function
\[
f^\star : \mathcal{X} \to \mathcal{Y}
\]
that governs the data:
\[
y = f^\star(x) + \text{noise}.
\]

Since $f^\star$ is unknown, we seek an approximation $f_\theta$ such that
\[
f_\theta \approx f^\star.
\]

\medskip

\noindent
To quantify this approximation, we introduce a \emph{loss function}
\[
\ell : \mathcal{Y} \times \mathcal{Y} \to \mathbb{R}.
\]

The quality of $f$ is measured by the \emph{risk}
\[
R(f) = \mathbb{E}_{(x,y)} \big[ \ell(f(x), y) \big].
\]

Since the data distribution is unknown, we instead minimize the \emph{empirical risk}
\[
\hat{R}_n(f) = \frac{1}{n} \sum_{i=1}^n \ell(f(x_i), y_i).
\]

\medskip

\noindent
\textbf{Learning problem.}
\[
\min_{\theta \in \Theta} \; \hat{R}_n(f_\theta).
\]

\medskip

This simple optimization problem is the foundation of nearly all modern machine learning algorithms.

\section{Examples from Regression and Classification}

\subsection{Regression}

In regression, the output space is $\mathcal{Y} = \mathbb{R}$.

A common choice of loss is the squared error:
\[
\ell(f(x), y) = (f(x) - y)^2.
\]

The empirical risk becomes
\[
\hat{R}_n(f) = \frac{1}{n} \sum_{i=1}^n (f(x_i) - y_i)^2.
\]

\medskip

\noindent
\textbf{Interpretation.}
\begin{itemize}
    \item Penalizes large deviations strongly
    \item Leads to analytically tractable solutions for linear models
\end{itemize}

\subsection{Classification}

In classification, $\mathcal{Y} = \{1, \dots, K\}$.

Instead of predicting labels directly, models often output scores or probabilities:
\[
f(x) \in \mathbb{R}^K.
\]

A common loss is the cross-entropy loss:
\[
\ell(f(x), y) = -\log p_\theta(y \mid x).
\]

\medskip

\noindent
\textbf{Interpretation.}
\begin{itemize}
    \item Encourages correct class probabilities
    \item Strong connection to maximum likelihood estimation
\end{itemize}

\section{From Approximation to Generalization}

A central issue arises: minimizing empirical risk does not guarantee good performance on new data.

\medskip

\noindent
\textbf{Key question.}  
Why should a function that fits the training data well also perform well on unseen data?

\medskip

This question leads to the study of:
\begin{itemize}
    \item generalization,
    \item model complexity,
    \item regularization,
    \item inductive bias.
\end{itemize}

\medskip

These ideas form the theoretical backbone of machine learning and will be developed throughout this book.

\section{A Unifying Perspective}

We summarize the learning process:

\begin{enumerate}
    \item Choose a model class $\mathcal{H}$
    \item Define a loss function $\ell$
    \item Solve an optimization problem
    \item Evaluate on unseen data
\end{enumerate}

\medskip

\noindent
\textbf{Three interacting components:}
\begin{itemize}
    \item \textbf{Approximation} (model expressivity)
    \item \textbf{Optimization} (ability to find good parameters)
    \item \textbf{Generalization} (performance on new data)
\end{itemize}

\medskip

Modern artificial intelligence, particularly deep learning, succeeds because it balances these three aspects at scale.

\section{Outlook}

This chapter introduced the central viewpoint of this book: learning as function approximation under uncertainty.

\medskip

In the next chapter, we develop the probabilistic framework underlying learning, connecting empirical risk minimization to statistical inference.

\medskip

\noindent
\textbf{Guiding principle.}  
Throughout this text, we will emphasize:
\begin{itemize}
    \item intuition grounded in mathematics,
    \item models as structured approximations,
    \item and algorithms as computational realizations of these ideas.
\end{itemize}
